% ---------------------------------------------------------------
% Baybayani User Evaluation Instrument
% Quality in use, based on ISO/IEC 25019:2023
% Respondents: customers and farmer-group administrators who were
% not part of the system's development.
% Build: pdflatex user-instrument-iso25019.tex (run twice)
% ---------------------------------------------------------------
\documentclass[10pt]{article}
% ---------------------------------------------------------------
% Study details shared by both evaluation instruments.
% Edit here once; both instruments pick the changes up on rebuild.
% Source: VSU appointment letters for Project DCST.0824-0825
% (September 10, 2024).
% ---------------------------------------------------------------

\newcommand{\SystemName}{Baybayani}
\newcommand{\ProjectTitle}{Farm to Table: A Management Information System for
  Farmers and Food Business Establishments in Baybay City, Leyte}
\newcommand{\ProjectCode}{DCST.0824-0825}

\newcommand{\ProjectLeader}{Mr.~Nelson H.~Tejara}
% Study/Component Leaders, alphabetical by surname
\newcommand{\ComponentLeaders}{%
  Mr.~Jomari Joseph A.~Barrera,
  Dr.~Jonah Flor O.~Maaghop,
  Mr.~Rodney M.~Maniego Jr.,
  Dr.~Jude B.~Rola, and
  Dr.~Magdalene C.~Unajan}

\newcommand{\University}{Visayas State University}
\newcommand{\Faculty}{Faculty of Computing}
\newcommand{\Department}{Department of Computer Science and Technology}
\newcommand{\Address}{Visca, Baybay City, Leyte 6521}

\usepackage{instrument}

\renewcommand{\InstrumentShortName}{User Evaluation (ISO/IEC 25019:2023)}
\renewcommand{\StatementHeader}{Statement \textit{/ Pahayag}}

% Cebuano line under an English line
\newcommand{\ceb}[1]{{\itshape\color{black!75}#1}}

\begin{document}

\InstrumentTitle{User Evaluation Instrument}%
  {Quality in Use of the \SystemName{} System, based on ISO/IEC 25019:2023\\
   \ceb{Pagtimbang-timbang sa Kalidad sa Paggamit sa \SystemName}}

% ===============================================================
\PartHeading{Part I. Informed Consent and Data Privacy Notice\\
  \ceb{\normalsize Bahin I. Pagtugot ug Pahibalo bahin sa Pribasiya sa Datos}}

Faculty researchers of the \Department{} of \University{} are evaluating
\SystemName, the farm-to-table ordering system of the research project
named above. We are asking you, as a user of the system, to rate how well it
serves you in your actual use.

\ceb{Ang mga magtutudlo-tigdukiduki sa \Department{} sa \University{} nagtimbang-timbang
sa \SystemName, ang farm-to-table nga sistema sa pag-order sa maong research
project. Gihangyo ka namo, isip usa ka naggamit sa sistema, nga hatagan og
grado kung unsa ka maayo ang pagserbisyo niini kanimo sa imong tinuod nga paggamit.}

\begin{itemize}
  \item \textbf{Voluntary.} Taking part is voluntary. You may skip any question or
    stop at any time without any effect on your use of \SystemName.\\
    \ceb{Boluntaryo ang pag-apil. Mahimo nimong lakawan ang bisan unsang pangutana o
    mohunong bisan kanus-a nga walay epekto sa imong paggamit sa \SystemName.}
  \item \textbf{Time.} Answering takes about 15 to 20 minutes.\\
    \ceb{Mga 15 ngadto sa 20 ka minuto ang pagtubag.}
  \item \textbf{Confidentiality.} Your name appears only on this page. This page is
    filed separately from your answers, which are identified only by a respondent
    number. Results are reported only as group summaries.\\
    \ceb{Kining panid lang ang adunay imong ngalan. Ibulag kini sa imong mga tubag,
    nga mailhan lang pinaagi sa respondent number. Ang mga resulta ireport lang
    isip katingban sa tanan.}
  \item \textbf{Data privacy.} Your information is collected, stored and used only for
    this evaluation, in line with the Data Privacy Act of 2012 (Republic Act
    No.~10173). Only the researchers can access it, and it is disposed of securely
    once the project is completed.\\
    \ceb{Ang imong impormasyon kolektahon, tipigan, ug gamiton lang alang niining
    pagtimbang-timbang, subay sa Data Privacy Act of 2012 (Republic Act No.~10173).
    Ang mga tigdukiduki lang ang maka-access niini, ug luwas kining labayon
    human mahuman ang proyekto.}
  \item \textbf{Eligibility.} Respondents must be 18 years old or older and must not
    have taken part in developing \SystemName.\\
    \ceb{Ang mga motubag kinahanglan 18 anyos pataas ug wala miapil sa paghimo sa \SystemName.}
\end{itemize}

For questions about this evaluation, approach any of the researchers at the
\Department, \University.\\
\ceb{Alang sa mga pangutana bahin niining pagtimbang-timbang, duol sa bisan kinsa sa
mga tigdukiduki sa \Department, \University.}

\vspace{4pt}
\opt{} I have read and understood this notice, and I agree to take part.\\
\hspace*{2.6ex}\ceb{Nabasa ug nasabtan nako kini nga pahibalo, ug mouyon ko nga moapil.}

\vspace{0.9cm}
\noindent
\begin{tabular}{@{}p{8.6cm}@{\hspace{1cm}}p{5cm}@{}}
  \rule{\linewidth}{0.4pt} & \rule{\linewidth}{0.4pt}\\
  \small Signature over printed name \ceb{/ Pirma ibabaw sa ngalan} &
  \small Date \ceb{/ Petsa}
\end{tabular}

\clearpage
% ===============================================================
\PartHeading{Part II. Respondent Profile\\
  \ceb{\normalsize Bahin II. Profile sa Motubag}}

Tick (\checkmark) the box that applies. \ceb{Butangi og tsek (\checkmark) ang kahon nga angay.}

{\small
\renewcommand{\arraystretch}{1.55}
\begin{tabularx}{\linewidth}{|>{\raggedright\arraybackslash}p{5.1cm}|X|}
\hline
\textbf{1. Role in \SystemName}\newline\ceb{Papel sa \SystemName} &
  \opt{Customer \ceb{(Kustomer)}}
  \opt{Administrator / farmer-group staff \ceb{(Administrador / kawani sa grupo sa mag-uuma)}}\\ \hline
\textbf{2. Type of customer}\newline\ceb{Klase sa kustomer}\newline{\footnotesize (customers only \ceb{/ kustomer lang})} &
  \opt{Individual or household \ceb{(Indibidwal o panimalay)}}
  \opt{Food business establishment, e.g.\ restaurant, eatery, carinderia \ceb{(Negosyo sa pagkaon)}}
  \opt{Reseller or retailer \ceb{(Tigbaligya pag-usab)}}
  \opt{Other \ceb{(Uban pa)}: \blank[2.6cm]}\\ \hline
\textbf{3. Age}\newline\ceb{Edad} &
  \opt{18--24} \opt{25--34} \opt{35--44} \opt{45--54} \opt{55--64} \opt{65 and above \ceb{(pataas)}}\\ \hline
\textbf{4. Sex}\newline\ceb{Sekso} &
  \opt{Male \ceb{(Lalaki)}} \opt{Female \ceb{(Babaye)}} \opt{Prefer not to say \ceb{(Dili gusto mosulti)}}\\ \hline
\textbf{5. How you use \SystemName}\newline\ceb{Unsaon nimo paggamit ang \SystemName}\newline{\footnotesize (tick all that apply \ceb{/ tsek tanan nga angay})} &
  \opt{Android app} \opt{Web browser on a phone \ceb{(browser sa cellphone)}}
  \opt{Web browser on a computer \ceb{(browser sa computer)}}\\ \hline
\textbf{6. How long you have used it}\newline\ceb{Unsa na kadugay nimo kini gigamit} &
  \opt{Less than 1 month \ceb{(Ubos sa 1 ka bulan)}} \opt{1--3 months \ceb{(bulan)}}
  \opt{4--6 months \ceb{(bulan)}} \opt{More than 6 months \ceb{(Labaw sa 6 ka bulan)}}\\ \hline
\textbf{7. How often you use it}\newline\ceb{Unsa ka sagad nimo kini gamiton} &
  \opt{Daily \ceb{(Adlaw-adlaw)}} \opt{A few times a week \ceb{(Pipila ka higayon sa usa ka semana)}}
  \opt{Weekly \ceb{(Kada semana)}} \opt{Less than weekly \ceb{(Dili kada semana)}}\\ \hline
\textbf{8. Online shopping or ordering before \SystemName}\newline\ceb{Online nga pagpalit o pag-order sa wala pa ang \SystemName} &
  \opt{Never \ceb{(Wala gyud)}} \opt{Sometimes \ceb{(Usahay)}} \opt{Often \ceb{(Kanunay)}}\\ \hline
\textbf{9. What you have done in \SystemName}\newline\ceb{Unsa na ang imong nabuhat sa \SystemName}\newline{\footnotesize (tick all that apply \ceb{/ tsek tanan nga angay})} &
  \renewcommand{\arraystretch}{1.15}%
  \begin{tabular}[t]{@{}p{0.47\linewidth}@{\hspace{0.03\linewidth}}p{0.47\linewidth}@{}}
  \multicolumn{2}{@{}l}{\textit{Customers \ceb{/ Kustomer}:}}\\
  \opt{Signed up \ceb{(Nag-sign up)}} & \opt{Browsed or searched products \ceb{(Nangita og produkto)}}\\
  \opt{Used the cart \ceb{(Migamit sa cart)}} & \opt{Placed an order \ceb{(Nag-order)}}\\
  \opt{Tracked an order \ceb{(Nagsubay sa order)}} & \opt{Cancelled an order \ceb{(Nagkansela og order)}}\\
  \opt{Chatted with the store \ceb{(Nakig-chat sa tindahan)}} & \opt{Read announcements \ceb{(Nagbasa sa pahibalo)}}\\
  \opt{Updated my profile \ceb{(Nag-update sa profile)}} & \\[3pt]
  \multicolumn{2}{@{}l}{\textit{Administrators \ceb{/ Administrador}:}}\\
  \opt{Managed products and variants \ceb{(Nagdumala sa produkto ug variant)}} &
  \opt{Adjusted or recorded stock \ceb{(Nag-adjust o nagrekord sa stock)}}\\
  \opt{Processed orders \ceb{(Nagproseso sa order)}} &
  \opt{Entered walk-in orders \ceb{(Nag-encode og walk-in nga order)}}\\
  \opt{Approved user accounts \ceb{(Nag-aprobar og account)}} &
  \opt{Posted announcements \ceb{(Nag-post og pahibalo)}}\\
  \opt{Changed store hours \ceb{(Nag-usab sa oras sa tindahan)}} &
  \opt{Replied to chats \ceb{(Mitubag sa chat)}}\\
  \opt{Viewed the dashboard or exported reports \ceb{(Mitan-aw sa dashboard o nag-export og report)}} & \\
  \end{tabular}\\ \hline
\end{tabularx}}

% ===============================================================
\PartHeading{Part III. Quality-in-Use Evaluation\\
  \ceb{\normalsize Bahin III. Pagtimbang-timbang sa Kalidad sa Paggamit}}

\textbf{Directions.} Read each statement and shade or tick (\checkmark) the circle
that best matches your own experience of using \SystemName. There are no right or
wrong answers. Answer Part~III-A, then the part for your role: Part~III-B if you
are a customer, Part~III-C if you are an administrator.\\
\ceb{\textbf{Giya.} Basaha ang matag pahayag ug lim-i o tsek-i (\checkmark) ang lingin nga
labing mohaum sa imong kaugalingong kasinatian sa paggamit sa \SystemName. Walay
husto o sayop nga tubag. Tubaga ang Bahin~III-A, dayon ang bahin para sa imong
papel: Bahin~III-B kung kustomer ka, Bahin~III-C kung administrador ka.}

\begin{center}
\begin{minipage}{\linewidth}
\centering
\small
\renewcommand{\arraystretch}{1.3}
\begin{tabular}{|c|c|l|l|c|}
\hline
\rowcolor{headshade}
\textbf{Rating} & \textbf{Code} & \textbf{Description} & \ceb{\textbf{Kahulugan}} & \textbf{Mean range}\\ \hline
5 & SA & Strongly Agree    & \ceb{Uyon kaayo}     & 4.21--5.00\\ \hline
4 & A  & Agree             & \ceb{Uyon}           & 3.41--4.20\\ \hline
3 & N  & Neutral           & \ceb{Dili sigurado}  & 2.61--3.40\\ \hline
2 & D  & Disagree          & \ceb{Dili uyon}      & 1.81--2.60\\ \hline
1 & SD & Strongly Disagree & \ceb{Dili gyud uyon} & 1.00--1.80\\ \hline
\end{tabular}\\[2pt]
{\footnotesize The mean range is used by the researchers to interpret the results.
\ceb{Gamiton sa mga tigdukiduki ang mean range sa paghubad sa mga resulta.}}
\end{minipage}
\end{center}

% ---------------------------------------------------------------
\PartHeading{Part III-A. For All Respondents \ceb{\normalsize / Para sa Tanang Motubag}}
\StartNumbering{A}

\Characteristic{1. Beneficialness \ceb{/ Kapuslanan}}{How much you benefit from
  using \SystemName: whether you can do what you need to do with it, whether people
  with different needs can use it, and whether it fits the work it is used for.
  \ceb{Unsa ka dako ang kaayohan nga imong makuha sa paggamit sa \SystemName.}}
\begin{ratingtable}
\SubChar{1.1 Usability \ceb{/ Kasayon sa paggamit}}
\BiItem{I can complete what I need to do in \SystemName{} without help from others.}
       {Mahimo nako ang akong kinahanglan buhaton sa \SystemName{} bisan walay tabang sa uban.}
\BiItem{I can finish my tasks in \SystemName{} quickly and in a few steps.}
       {Dali ra nako mahuman ang akong mga buluhaton sa \SystemName, sa pipila lang ka lakang.}
\BiItem{The results I get are complete and correct (e.g., prices, quantities, order details).}
       {Kompleto ug husto ang mga resulta nga akong makuha (pananglitan, presyo, gidaghanon, detalye sa order).}
\BiItem{The screens, labels and buttons are easy to understand.}
       {Sayon sabton ang mga screen, label, ug button.}
\BiItem{I am satisfied with using \SystemName.}
       {Kontento ko sa paggamit sa \SystemName.}
\SubChar{1.2 Accessibility \ceb{/ Kaabot sa tanan}}
\BiItem{The text, pictures and buttons are easy to read and see on my device.}
       {Sayon basahon ug tan-awon ang mga sulat, hulagway, ug button sa akong device.}
\BiItem{I can use \SystemName{} well on the device I have (phone app, phone browser or computer).}
       {Magamit nako og maayo ang \SystemName{} sa device nga naa nako (app sa cellphone, browser sa cellphone, o computer).}
\BiItem{\SystemName{} can be used even by people who are not used to apps or online shopping.}
       {Magamit ang \SystemName{} bisan sa mga tawo nga dili anad sa mga app o online shopping.}
\BiItem{I can set the display (e.g., light or dark mode) to what is comfortable for me.}
       {Mausab nako ang hitsura (pananglitan, light o dark mode) sumala sa komportable para nako.}
\SubChar{1.3 Suitability \ceb{/ Kaangayan}}
\BiItem{\SystemName{} fits the way the farmer group sells and the way buyers buy produce.}
       {Mohaum ang \SystemName{} sa paagi sa pagbaligya sa grupo sa mag-uuma ug sa paagi sa pagpalit sa mga mamalitay.}
\BiItem{\SystemName{} works whenever I need it.}
       {Molihok ang \SystemName{} matag higayon nga akong kinahanglanon.}
\BiItem{\SystemName{} responds fast enough, even on mobile data.}
       {Paspas igo ang tubag sa \SystemName, bisan sa mobile data.}
\BiItem{Notifications and chat messages arrive on time.}
       {Moabot sa saktong oras ang mga notification ug mensahe sa chat.}
\end{ratingtable}
\CommentBox[2]{Comments on beneficialness \ceb{/ Komento bahin sa kapuslanan}:}

\Characteristic{2. Freedom from Risk \ceb{/ Kagawasan gikan sa Peligro}}{How well
  \SystemName{} keeps you away from financial loss, harm to your health and safety,
  and harm to the community and the environment.
  \ceb{Unsa ka maayo nga gilikay ka sa \SystemName{} gikan sa kapildihan, kadaot sa
  panglawas ug kaluwasan, ug kadaot sa komunidad ug kinaiyahan.}}
\begin{ratingtable}
\SubChar{2.1 Freedom from economic risk \ceb{/ Kagawasan gikan sa peligro sa kwarta}}
\BiItem{The prices and amounts I see are the ones actually charged.}
       {Ang presyo ug kantidad nga akong makita mao gyud ang sukton.}
\BiItem{\SystemName{} helps me avoid losses from wrong prices, wrong quantities or lost orders.}
       {Makatabang ang \SystemName{} nga malikayan nako ang kapildihan tungod sa sayop nga presyo, sayop nga gidaghanon, o nawala nga order.}
\BiItem{Using \SystemName{} saves me time, fare or other costs.}
       {Makadaginot ko og oras, pamasahe, o ubang gasto tungod sa \SystemName.}
\SubChar{2.2 Freedom from health risk \ceb{/ Kagawasan gikan sa peligro sa panglawas}}
\BiItem{Product information (photos, descriptions, availability) helps me get fresh, safe produce.}
       {Ang impormasyon sa produkto (hulagway, deskripsyon, kung naa pa) makatabang nako nga makakuha og presko ug luwas nga produkto.}
\BiItem{Using \SystemName{} does not cause me stress, eye strain or discomfort.}
       {Ang paggamit sa \SystemName{} dili makahatag nako og stress, kakapoy sa mata, o kahasol.}
\SubChar{2.3 Freedom from human life risk \ceb{/ Kagawasan gikan sa peligro sa kinabuhi}}
\BiItem{\SystemName{} reduces my need to travel or wait in crowded places to buy or sell.}
       {Makunhuran sa \SystemName{} ang akong panginahanglan nga mobiyahe o maghulat sa daghang tawo aron mopalit o mobaligya.}
\BiItem{\SystemName{} never requires me to do anything that puts my personal safety at risk.}
       {Wala gyud ko pugsa sa \SystemName{} nga mobuhat og bisan unsa nga makabutang sa akong kaluwasan sa peligro.}
\SubChar{2.4 Freedom from environmental and societal risk \ceb{/ Kagawasan gikan sa peligro sa kinaiyahan ug katilingban}}
\BiItem{\SystemName{} helps reduce wasted produce (e.g., unsold or spoiled goods).}
       {Makatabang ang \SystemName{} sa pagkunhod sa nausik nga produkto (pananglitan, wala mabaligya o nadunot).}
\BiItem{\SystemName{} supports local farmers and the community of Baybay City.}
       {Gisuportahan sa \SystemName{} ang mga lokal nga mag-uuma ug ang komunidad sa Dakbayan sa Baybay.}
\BiItem{\SystemName{} reduces paper use (e.g., handwritten orders, receipts and lists).}
       {Makunhuran sa \SystemName{} ang paggamit og papel (pananglitan, sinulat nga order, resibo, ug lista).}
\end{ratingtable}
\CommentBox[2]{Comments on freedom from risk \ceb{/ Komento bahin sa kagawasan gikan sa peligro}:}

\Characteristic{3. Acceptability \ceb{/ Kadawaton}}{How favourably you respond to
  \SystemName: your experience with it over time, how far you trust it, and whether
  it follows the rules that apply to it.
  \ceb{Unsa ka maayo ang imong pagdawat sa \SystemName: ang imong kasinatian niini,
  ang imong pagsalig niini, ug kung nagsunod kini sa mga lagda.}}
\begin{ratingtable}
\SubChar{3.1 Experience \ceb{/ Kasinatian}}
\BiItem{My experience with \SystemName{} has been positive.}
       {Maayo ang akong kasinatian sa \SystemName.}
\BiItem{The more I use \SystemName, the easier it becomes.}
       {Kon mas kanunay nako gamiton ang \SystemName, mas mosayon kini.}
\BiItem{I will continue using \SystemName.}
       {Padayon nakong gamiton ang \SystemName.}
\BiItem{I would recommend \SystemName{} to others.}
       {Irekomendar nako ang \SystemName{} sa uban.}
\SubChar{3.2 Trustworthiness \ceb{/ Kasaligan}}
\BiItem{I trust that \SystemName{} keeps my personal information, including my valid ID, safe.}
       {Nagsalig ko nga luwas sa \SystemName{} ang akong personal nga impormasyon, lakip ang akong valid ID.}
\BiItem{I trust that the orders, stock and records in \SystemName{} are accurate.}
       {Nagsalig ko nga husto ang mga order, stock, ug rekord sa \SystemName.}
\BiItem{\SystemName{} behaves the way I expect it to.}
       {Molihok ang \SystemName{} sumala sa akong gipaabot.}
\SubChar{3.3 Compliance \ceb{/ Pagsunod sa lagda}}
\BiItem{I am confident that \SystemName{} handles personal data in line with the Data Privacy Act of 2012 (RA 10173).}
       {Kumpiyansa ko nga gidumala sa \SystemName{} ang personal nga datos subay sa Data Privacy Act of 2012 (RA 10173).}
\BiItem{\SystemName{} follows the farmer group's rules (e.g., store hours, account approval, order cancellation).}
       {Gisunod sa \SystemName{} ang mga lagda sa grupo sa mag-uuma (pananglitan, oras sa tindahan, pag-aprobar sa account, pagkansela sa order).}
\end{ratingtable}
\CommentBox[2]{Comments on acceptability \ceb{/ Komento bahin sa kadawaton}:}

% ---------------------------------------------------------------
\PartHeading{Part III-B. For Customers Only \ceb{\normalsize / Para sa Kustomer Lang}}
\StartNumbering{B}
\begin{ratingtable}
\SubChar{1. Beneficialness: usability and suitability \ceb{/ Kapuslanan: kasayon sa paggamit ug kaangayan}}
\BiItem{I can easily find products using categories, tags and search.}
       {Dali nakong makit-an ang mga produkto pinaagi sa kategorya, tag, ug search.}
\BiItem{I can easily choose a product variant and quantity and add it to my cart.}
       {Dali ra nako mapili ang variant ug gidaghanon sa produkto ug idugang kini sa akong cart.}
\BiItem{I can easily review, edit and check out my cart.}
       {Dali ra nako masusi, ma-edit, ug ma-check out ang akong cart.}
\BiItem{I can easily follow the status of my order (Pending, Ready, Completed or Cancelled).}
       {Dali ra nako masubay ang status sa akong order (Pending, Ready, Completed, o Cancelled).}
\BiItem{Messaging the store through the chat is easy, and my concerns get answered.}
       {Sayon ang pag-message sa tindahan pinaagi sa chat, ug matubag ang akong mga pangutana.}
\BiItem{Announcements and store hours tell me what I can buy and when.}
       {Ang mga pahibalo ug oras sa tindahan nagpahibalo nako kung unsa ang akong mapalit ug kanus-a.}
\BiItem{Signing up and getting my account approved was straightforward.}
       {Dili lisod ang pag-sign up ug ang pag-aprobar sa akong account.}
\SubChar{2. Freedom from economic risk \ceb{/ Kagawasan gikan sa peligro sa kwarta}}
\BiItem{Retail and wholesale prices are shown clearly before I order.}
       {Klaro nga gipakita ang retail ug wholesale nga presyo sa dili pa ko mag-order.}
\BiItem{I am told clearly when an item is out of stock or my order is cancelled, and why.}
       {Klaro nga ipahibalo nako kung wala nay stock ang usa ka produkto o kung gikansela ang akong order, ug ngano.}
\SubChar{3. Acceptability: experience \ceb{/ Kadawaton: kasinatian}}
\BiItem{I prefer ordering through \SystemName{} to ordering in person, by text or by phone call.}
       {Mas gusto nako ang pag-order pinaagi sa \SystemName{} kaysa personal, text, o tawag.}
\end{ratingtable}
\CommentBox[2]{Comments on ordering as a customer \ceb{/ Komento bahin sa pag-order isip kustomer}:}

% ---------------------------------------------------------------
\PartHeading{Part III-C. For Administrators Only \ceb{\normalsize / Para sa Administrador Lang}}
\StartNumbering{C}
\begin{ratingtable}
\SubChar{1. Beneficialness: usability and suitability \ceb{/ Kapuslanan: kasayon sa paggamit ug kaangayan}}
\BiItem{Adding and editing products, variants, prices and photos is easy.}
       {Sayon ang pagdugang ug pag-edit sa mga produkto, variant, presyo, ug hulagway.}
\BiItem{Processing orders, including walk-in orders, is quick and easy.}
       {Paspas ug sayon ang pagproseso sa mga order, lakip ang mga walk-in nga order.}
\BiItem{Reviewing and approving customer accounts, including their valid IDs, is easy.}
       {Sayon ang pagsusi ug pag-aprobar sa mga account sa kustomer, lakip ang ilang valid ID.}
\BiItem{Posting announcements and changing store hours is easy.}
       {Sayon ang pag-post og pahibalo ug pag-usab sa oras sa tindahan.}
\BiItem{The dashboard gives me the sales, order and customer information I need to make decisions.}
       {Gihatag sa dashboard ang impormasyon bahin sa halin, order, ug kustomer nga akong kinahanglan sa paghimo og desisyon.}
\BiItem{The stock records (acquisitions, sales, losses, restocks from cancellations) match the actual inventory.}
       {Mohaum sa tinuod nga inventory ang mga rekord sa stock (acquisition, sale, loss, restock gikan sa kanselado).}
\BiItem{Low-stock notifications help me restock in time.}
       {Makatabang ang mga notification sa gamay na nga stock aron makapuno ko sa saktong panahon.}
\SubChar{2. Freedom from economic risk \ceb{/ Kagawasan gikan sa peligro sa kwarta}}
\BiItem{\SystemName{} helps the farmer group reduce losses from overselling, spoilage or recording errors.}
       {Makatabang ang \SystemName{} sa grupo sa mag-uuma sa pagkunhod sa kapildihan tungod sa sobra nga pagbaligya, pagkadunot, o sayop sa pagrekord.}
\BiItem{\SystemName{} warns me before actions that cannot easily be undone (e.g., deleting, cancelling, adjusting stock).}
       {Gipasidan-an ko sa \SystemName{} sa dili pa ko mobuhat og aksyon nga dili dali mabawi (pananglitan, pagtangtang, pagkansela, pag-adjust sa stock).}
\SubChar{3. Acceptability \ceb{/ Kadawaton}}
\BiItem{Only authorized administrators can open the admin functions.}
       {Ang mga awtorisado lang nga administrador ang makaabli sa mga admin function.}
\BiItem{The reports exported from \SystemName{} are acceptable for the farmer group's records and reporting.}
       {Madawat para sa mga rekord ug report sa grupo sa mag-uuma ang mga report nga gi-export gikan sa \SystemName.}
\BiItem{\SystemName{} has made managing the farmer group's sales easier than before.}
       {Mas nasayon ang pagdumala sa halin sa grupo sa mag-uuma tungod sa \SystemName{} kaysa kaniadto.}
\end{ratingtable}
\CommentBox[2]{Comments on managing the store as an administrator \ceb{/ Komento bahin sa pagdumala sa tindahan isip administrador}:}

% ===============================================================
% Part IV is kept on one page.
\par\noindent\begin{minipage}{\linewidth}
\PartHeading{Part IV. Comments and Suggestions\\
  \ceb{\normalsize Bahin IV. Mga Komento ug Sugyot}}

\CommentBox[3]{1. What did you like most about \SystemName?\\
  \ceb{Unsa ang labing nimo nagustohan sa \SystemName?}}
\vspace{4pt}
\CommentBox[3]{2. What problems did you encounter while using it?\\
  \ceb{Unsa nga mga problema ang imong nasugatan samtang gigamit kini?}}
\vspace{4pt}
\CommentBox[3]{3. What would you suggest to improve \SystemName?\\
  \ceb{Unsa ang imong isugyot aron mapaayo ang \SystemName?}}

\vspace{10pt}
\begin{center}
\textbf{Thank you for your time and honest answers.}\\
\ceb{Daghang salamat sa imong panahon ug matinuoron nga mga tubag.}
\end{center}
\end{minipage}

\end{document}
